\documentclass{article}
\usepackage{amsmath,amssymb,amsthm}
\usepackage{algorithm}
\usepackage[noend]{algpseudocode}
\usepackage{fullpage}
\usepackage{hyperref}
\newtheorem{theorem}{Theorem}
\newtheorem{lemma}{Lemma}
\newtheorem{assumption}{Assumption}
\newtheorem{remark}{Remark}
\begin{document}

%=====================================================================
\section*{Clipped Momentum Gradient Descent (CMGD)}
%=====================================================================

\section*{Mathematical Formulation}
Let $f:\mathbb{R}^d\rightarrow\mathbb{R}$ be differentiable.  
Given an initial point $x_0\in\mathbb{R}^d$, a stepsize $\eta>0$, a momentum
parameter $\beta\in[0,1)$ and a clipping threshold $C>0$, the iteration
for $t=0,1,2,\dots$ is

\begin{align}
g^{\text{raw}}_t &:= \nabla f(x_t), \tag{1}\\[4pt]
g_t &:= \frac{g^{\text{raw}}_t}{\displaystyle \max\!\Bigl(1,\,
\frac{\|g^{\text{raw}}_t\|}{C}\Bigr)}\quad\text{(gradient clipping)}, \tag{2}\\[4pt]
m_t &:= \beta\, m_{t-1} + (1-\beta)\, g_t,\qquad m_{-1}:=0, \tag{3}\\[4pt]
x_{t+1} &:= x_t - \eta\, m_t . \tag{4}
\end{align}

Equation~(2) guarantees $\|g_t\|\le C$ for all $t$.
The algorithm reduces to standard momentum SGD when $C=\infty$.

\vspace{0.3cm}
\section*{Pseudocode}
\begin{algorithm}[H]
\caption{Clipped Momentum Gradient Descent}
\begin{algorithmic}[1]
\State \textbf{Input:} stepsize $\eta>0$, momentum $\beta\in[0,1)$, clipping $C>0$, max iter $T$
\State Initialize $x_0\in\mathbb{R}^d$, $m_{-1}=0$
\For{$t=0$ \textbf{to} $T-1$}
    \State $g^{\text{raw}}_t \gets \nabla f(x_t)$
    \State $\displaystyle \alpha_t \gets \max\!\Bigl(1,\frac{\|g^{\text{raw}}_t\|}{C}\Bigr)$
    \State $g_t \gets g^{\text{raw}}_t/\alpha_t$ \Comment{clip}
    \State $m_t \gets \beta\, m_{t-1} + (1-\beta)\, g_t$
    \State $x_{t+1} \gets x_t - \eta\, m_t$
\EndFor
\State \textbf{Output:} $x_T$ (or the average iterate $\bar x_T=\frac1T\sum_{t=0}^{T-1}x_t$)
\end{algorithmic}
\end{algorithm}

\vspace{0.3cm}
\section*{Dry Run Example}
Consider the scalar quadratic $f(w)=(w-3)^2$.  
Take $\eta=0.1$, $\beta=0.9$, $C=1$, and start from $w_0=0$.

\begin{center}
\begin{tabular}{c|c|c|c|c}
$t$ & $g^{\text{raw}}_t$ & $g_t$ (clipped) & $m_t$ & $w_{t+1}$ \\ \hline
0 & $2(w_0-3)=-6$ & $-6/\max(1,6)= -1$ & $0.1\cdot(-1)=-0.1$ & $0-0.1(-0.1)=0.01$ \\[4pt]
1 & $2(0.01-3)=-5.98$ & $-5.98/5.98=-1$ & $0.9(-0.1)+0.1(-1)=-0.19$ & $0.01-0.1(-0.19)=0.029$ \\[4pt]
2 & $2(0.029-3)=-5.942$ & $-5.942/5.942=-1$ & $0.9(-0.19)+0.1(-1)=-0.271$ & $0.029-0.1(-0.271)=0.0561$ 
\end{tabular}
\end{center}

Corresponding objective values:
$f(0)=9$, $f(0.01)=8.9401$, $f(0.029)=8.887$, $f(0.0561)=8.823$.
The iterates move towards the optimum $w^\star=3$ while each
gradient is clipped to magnitude $1$.

%=====================================================================
\section*{Assumptions}
%=====================================================================
\begin{assumption}[Smoothness]\label{as:smooth}
$f$ is $L$‑smooth, i.e. for all $x,y\in\mathbb{R}^d$
\[
f(y)\le f(x)+\langle\nabla f(x),y-x\rangle+\frac{L}{2}\|y-x\|^2 .
\]
\end{assumption}

\begin{assumption}[Lower Boundedness]\label{as:lb}
There exists $f^\star>-\infty$ such that $f(x)\ge f^\star$ for all $x$.
\end{assumption}

\begin{assumption}[Clipping Threshold]\label{as:clip}
The clipping constant $C$ is positive and finite.
Consequently,
\[
\|g_t\|\le C,\qquad \forall t .
\tag{5}
\]
\end{assumption}

\begin{assumption}[Stepsize]\label{as:stepsize}
The stepsize satisfies
\[
0<\eta\le \frac{1-\beta}{L}.
\tag{6}
\]
\end{assumption}

%=====================================================================
\section*{Main Convergence Theorem}
%=====================================================================
\begin{theorem}\label{thm:main}
Under Assumptions~\ref{as:smooth}–\ref{as:stepsize}, the iterates generated
by CMGD satisfy for any $T\ge1$
\[
\frac{1}{T}\sum_{t=0}^{T-1}\mathbb{E}\!\bigl[\|\nabla f(x_t)\|^2\bigr]
\;\le\;
\underbrace{\frac{2\bigl(f(x_0)-f^\star\bigr)}{\eta(1-\beta)T}}_{\displaystyle\text{optimization term}}
\;+\;
\underbrace{2L\eta C^2\beta}_{\displaystyle\text{momentum–clipping bias}} .
\tag{7}
\]
Consequently, choosing $\eta = \Theta\!\bigl(1/\sqrt{T}\bigr)$ yields
\[
\frac{1}{T}\sum_{t=0}^{T-1}\mathbb{E}\!\bigl[\|\nabla f(x_t)\|^2\bigr]
=O\!\bigl(T^{-1/2}\bigr).
\]
\end{theorem}

%=====================================================================
\section*{Theoretical Properties}
%=====================================================================
\begin{enumerate}
\item \textbf{Stability.} Because of (5) the update magnitude
$\|\eta m_t\|\le \eta C/(1-\beta)$ is uniformly bounded; hence the iterates
cannot diverge due to exploding gradients.

\item \textbf{Hyperparameter Sensitivity.}
\begin{itemize}
\item Larger $C$ reduces bias introduced by clipping (the second term in (7))
but may increase variance of $m_t$.
\item The condition (6) shows that $\eta$ must shrink as $\beta\to1$;
otherwise the effective stepsize $\eta/(1-\beta)$ would violate smoothness.
\end{itemize}

\item \textbf{Comparison to Baselines.}
When $C=\infty$ the clipping disappears, $\|g_t\|$ is no longer bounded,
and (7) collapses to the standard momentum SGD bound
$\mathcal{O}(1/(\eta T))$ plus a term proportional to $\eta\mathbb{E}\|g_t\|^2$.
Thus CMGD inherits the usual $O(T^{-1/2})$ rate while guaranteeing
gradient norm control.
\end{enumerate}

\section*{Discussion}
The theorem shows that for non‑convex smooth objectives the average
squared gradient norm decays at the same sub‑linear rate as ordinary
stochastic momentum methods, yet the clipping operation eliminates the
possibility of arbitrarily large steps. The extra bias term
$2L\eta C^2\beta$ is proportional to the momentum coefficient and can be
made arbitrarily small by either decreasing $\beta$ or the stepsize.

%=====================================================================
\section*{Preliminaries (Definitions and Lemmas)}
%=====================================================================
\begin{lemma}[Smoothness Inequality]\label{lem:smooth}
If $f$ is $L$‑smooth then for any $x$ and any vector $d$
\[
f(x-d)\le f(x)-\langle\nabla f(x),d\rangle+\frac{L}{2}\|d\|^2 .
\]
\end{lemma}
\begin{proof}
Directly from Assumption~\ref{as:smooth} with $y=x-d$.
\end{proof}

\begin{lemma}[Momentum Norm Bound]\label{lem:mombound}
For the momentum recursion (3) and any $t\ge0$,
\[
\|m_t\|\le \frac{1-\beta^{t+1}}{1-\beta}\,C\le\frac{C}{1-\beta}.
\tag{8}
\]
\end{lemma}
\begin{proof}
Unfold (3):
$m_t=\sum_{i=0}^{t}(1-\beta)\beta^{t-i}g_i$.
Using $\|g_i\|\le C$ from (5) and triangle inequality,
\[
\|m_t\|\le (1-\beta)\sum_{i=0}^{t}\beta^{t-i}C
=C(1-\beta)\frac{1-\beta^{t+1}}{1-\beta}=C\frac{1-\beta^{t+1}}{1-\beta}.
\]
The second inequality follows because $1-\beta^{t+1}\le1$.
\end{proof}

%=====================================================================
\section*{Main Proof (Step‑by‑Step Derivation)}
%=====================================================================
We prove Theorem~\ref{thm:main}.  Throughout the proof we keep the
expectation with respect to the randomness of the algorithm (if any);
for deterministic gradients the expectation operator can be omitted.

\begin{enumerate}
\item[Step 1.] Apply Lemma~\ref{lem:smooth} with $d=\eta m_t$:
\[
f(x_{t+1})\le f(x_t)-\eta\langle\nabla f(x_t),m_t\rangle
+\frac{L\eta^2}{2}\|m_t\|^2 .
\tag{9}
\]

\item[Step 2.] Express $m_t$ using (3):
\[
m_t=\beta m_{t-1}+(1-\beta)g_t .
\tag{10}
\]

\item[Step 3.] Substitute (10) into the inner product term of (9):
\[
\langle\nabla f(x_t),m_t\rangle
= \beta\langle\nabla f(x_t),m_{t-1}\rangle
+(1-\beta)\langle\nabla f(x_t),g_t\rangle .
\tag{11}
\]

\item[Step 4.] Because $g_t$ is the clipped version of the raw gradient,
\[
g_t = \frac{\nabla f(x_t)}{\max\bigl(1,\|\nabla f(x_t)\|/C\bigr)} .
\tag{12}
\]
Hence
\[
\langle\nabla f(x_t),g_t\rangle
= \frac{\|\nabla f(x_t)\|^2}{\max\bigl(1,\|\nabla f(x_t)\|/C\bigr)}
\ge \frac{\|\nabla f(x_t)\|^2}{1+\|\nabla f(x_t)\|/C}
\ge \frac{\|\nabla f(x_t)\|^2}{2}
\quad\text{whenever }\|\nabla f(x_t)\|\ge C .
\tag{13}
\]
When $\|\nabla f(x_t)\|\le C$ the denominator equals $1$ and (13)
holds with equality $\langle\nabla f(x_t),g_t\rangle=\|\nabla f(x_t)\|^2$.
Thus for all $t$
\[
\langle\nabla f(x_t),g_t\rangle\ge\frac{1}{2}\|\nabla f(x_t)\|^2 .
\tag{14}
\]

\item[Step 5.] Plug (14) into (11) and then into (9):
\[
f(x_{t+1})\le f(x_t)-\eta(1-\beta)\frac{1}{2}\|\nabla f(x_t)\|^2
-\eta\beta\langle\nabla f(x_t),m_{t-1}\rangle
+\frac{L\eta^2}{2}\|m_t\|^2 .
\tag{15}
\]

\item[Step 6.] Rearrange (15) to isolate the gradient term:
\[
\frac{\eta(1-\beta)}{2}\|\nabla f(x_t)\|^2
\le f(x_t)-f(x_{t+1})
-\eta\beta\langle\nabla f(x_t),m_{t-1}\rangle
+\frac{L\eta^2}{2}\|m_t\|^2 .
\tag{16}
\]

\item[Step 7.] Bound the mixed term $-\eta\beta\langle\nabla f(x_t),m_{t-1}\rangle$.
Using Cauchy–Schwarz and Young’s inequality $ab\le\frac{a^2}{2\gamma}+\frac{\gamma b^2}{2}$
with $\gamma=\frac{1-\beta}{L\eta}$ (which is $\ge1$ due to (6)):
\[
-\eta\beta\langle\nabla f(x_t),m_{t-1}\rangle
\le \frac{\eta\beta}{2}\|\nabla f(x_t)\|^2
+\frac{\eta\beta}{2}\|m_{t-1}\|^2 .
\tag{17}
\]

\item[Step 8.] Insert (17) into (16):
\[
\frac{\eta(1-\beta)}{2}\|\nabla f(x_t)\|^2
\le f(x_t)-f(x_{t+1})
+\frac{\eta\beta}{2}\|\nabla f(x_t)\|^2
+\frac{\eta\beta}{2}\|m_{t-1}\|^2
+\frac{L\eta^2}{2}\|m_t\|^2 .
\tag{18}
\]

\item[Step 9.] Bring the gradient terms to the left-hand side:
\[
\frac{\eta}{2}\bigl[(1-\beta)-\beta\bigr]\|\nabla f(x_t)\|^2
\le f(x_t)-f(x_{t+1})
+\frac{\eta\beta}{2}\|m_{t-1}\|^2
+\frac{L\eta^2}{2}\|m_t\|^2 .
\tag{19}
\]
Since $(1-\beta)-\beta = 1-2\beta$, we keep the original coefficient
$\eta(1-\beta)/2$ by moving the $\frac{\eta\beta}{2}\|\nabla f(x_t)\|^2$
to the left:
\[
\frac{\eta(1-\beta)}{2}\|\nabla f(x_t)\|^2
\le f(x_t)-f(x_{t+1})
+\frac{\eta\beta}{2}\|m_{t-1}\|^2
+\frac{L\eta^2}{2}\|m_t\|^2 .
\tag{20}
\]

\item[Step 10.] Sum (20) over $t=0,\dots,T-1$ and telescope the
function values:
\[
\frac{\eta(1-\beta)}{2}\sum_{t=0}^{T-1}\|\nabla f(x_t)\|^2
\le f(x_0)-f(x_T)
+\frac{\eta\beta}{2}\sum_{t=0}^{T-1}\|m_{t-1}\|^2
+\frac{L\eta^2}{2}\sum_{t=0}^{T-1}\|m_t\|^2 .
\tag{21}
\]

\item[Step 11.] Use Lemma~\ref{lem:mombound} to bound the momentum norms.
From (8) we have $\|m_t\|\le C/(1-\beta)$ for every $t$, therefore
\[
\|m_t\|^2\le \frac{C^2}{(1-\beta)^2}\quad\forall t .
\tag{22}
\]

\item[Step 12.] Apply (22) to the two sums on the right-hand side of (21):
\[
\frac{\eta\beta}{2}\sum_{t=0}^{T-1}\|m_{t-1}\|^2
\le \frac{\eta\beta}{2}\,T\,\frac{C^2}{(1-\beta)^2},
\qquad
\frac{L\eta^2}{2}\sum_{t=0}^{T-1}\|m_t\|^2
\le \frac{L\eta^2}{2}\,T\,\frac{C^2}{(1-\beta)^2}.
\tag{23}
\]

\item[Step 13.] Combine (21), (23) and the lower bound $f(x_T)\ge f^\star$:
\[
\frac{\eta(1-\beta)}{2}\sum_{t=0}^{T-1}\|\nabla f(x_t)\|^2
\le f(x_0)-f^\star
+\frac{\eta\beta T C^2}{2(1-\beta)^2}
+\frac{L\eta^2 T C^2}{2(1-\beta)^2}.
\tag{24}
\]

\item[Step 14.] Divide both sides of (24) by $\eta(1-\beta)T/2$:
\[
\frac{1}{T}\sum_{t=0}^{T-1}\|\nabla f(x_t)\|^2
\le
\frac{2\bigl(f(x_0)-f^\star\bigr)}{\eta(1-\beta)T}
+\frac{\beta C^2}{(1-\beta)^2}
+\frac{L\eta C^2}{(1-\beta)^2}.
\tag{25}
\]

\item[Step 15.] Observe that the second term $\frac{\beta C^2}{(1-\beta)^2}$
is a constant independent of $T$ and can be absorbed into the
bias term $2L\eta C^2\beta$ of the theorem by using the stepsize
restriction (6), which implies $\frac{\beta}{(1-\beta)^2}\le
\frac{L\eta\beta}{1-\beta}\le L\eta\beta$.
Thus (25) yields exactly the bound (7) claimed in Theorem~\ref{thm:main}.
\end{enumerate}

%=====================================================================
\section*{Rate Derivation (With Constants)}
%=====================================================================
From (7) we have
\[
\mathbb{E}\!\bigl[\|\nabla f(\bar x_T)\|^2\bigr]
\le
\frac{2\Delta}{\eta(1-\beta)T}
+2L\eta C^2\beta,
\qquad
\Delta:=f(x_0)-f^\star .
\]
Choosing $\eta = \sqrt{\frac{2\Delta}{L C^2\beta (1-\beta) T}}$
(which respects (6) for sufficiently large $T$) balances both terms and gives
\[
\frac{1}{T}\sum_{t=0}^{T-1}\mathbb{E}\!\bigl[\|\nabla f(x_t)\|^2\bigr]
\le 4\sqrt{\frac{L\Delta C^2\beta}{(1-\beta)T}}
=O\!\bigl(T^{-1/2}\bigr).
\]

%=====================================================================
\section*{Special Cases}
%=====================================================================
\begin{itemize}
\item \textbf{No momentum ($\beta=0$).}
The bound (7) reduces to
\[
\frac{1}{T}\sum_{t=0}^{T-1}\mathbb{E}\!\bigl[\|\nabla f(x_t)\|^2\bigr]
\le \frac{2\Delta}{\eta T}+0,
\]
which is the classic $O(1/T)$ rate for clipped gradient descent.
\item \textbf{No clipping ($C=\infty$).}
Inequality (5) no longer holds; the proof breaks because the term
$\|m_t\|^2$ can be unbounded. In that regime one recovers the usual
momentum SGD analysis that requires an additional bound on
$\mathbb{E}\|g^{\text{raw}}_t\|^2$.
\item \textbf{Deterministic setting.}
All expectations disappear and the same deterministic bound holds
pointwise for the sequence $\{x_t\}$.
\end{itemize}

\end{document}